\documentclass[12pt,a4paper]{article}

% ---------------------------
% Pacotes essenciais
% ---------------------------
\usepackage[utf8]{inputenc}
\usepackage[T1]{fontenc}
\usepackage[brazil]{babel}
\usepackage{helvet}
\renewcommand{\familydefault}{\sfdefault}
\usepackage{geometry}
\usepackage{setspace}
\usepackage[table]{xcolor}
\usepackage{titlesec}
\usepackage{fancyhdr}
\usepackage{longtable}
\usepackage{tabularx}
\usepackage{array}
\usepackage{indentfirst}
\usepackage{ragged2e}
\usepackage{xurl}
\usepackage{hyperref}

% ---------------------------
% Paleta de Cores Profissional
% ---------------------------
\definecolor{AzulCorporativo}{RGB}{0, 82, 155}
\definecolor{AzulPetroleo}{RGB}{23, 104, 133}
\definecolor{CinzaEscuro}{RGB}{64, 64, 64}
\definecolor{LaranjaDestaque}{RGB}{230, 115, 0}

% ---------------------------
% Configuração de Links (Hyperref)
% ---------------------------
\hypersetup{
    colorlinks=true,
    linkcolor=AzulPetroleo,
    urlcolor=AzulCorporativo,
    citecolor=LaranjaDestaque,
    pdftitle={Análise de Requisitos},
    pdfauthor={Lucas Dias Noronha}
}

% ---------------------------
% Configuração de Títulos (Titlesec)
% ---------------------------
\titleformat{\section}
{\color{AzulCorporativo}\normalfont\Large\bfseries}
{\thesection}{1em}{}

\titleformat{\subsection}
{\color{AzulPetroleo}\normalfont\large\bfseries}
{\thesubsection}{1em}{}

% ---------------------------
% Configuração de página e Cabeçalho
% ---------------------------
\geometry{
    left=3cm,
    right=2cm,
    top=3cm,
    bottom=2cm,
    headheight=45pt
}

\onehalfspacing
\setlength{\emergencystretch}{2em}

\pagestyle{fancy}
\fancyhf{}
\renewcommand{\headrulewidth}{0pt}

\fancyhead[C]{%
    \colorbox{AzulCorporativo}{%
        \makebox[\dimexpr\textwidth-2\fboxsep\relax][l]{%
            \textcolor{white}{\textbf{\ \ Análise de Requisitos}}%
        }%
    }%
}

\fancyfoot[C]{\color{CinzaEscuro}\thepage}

% ---------------------------
% Variáveis do documento
% ---------------------------
\newcommand{\projeto}{Modelo Analítico de Dados para E-commerce}
\newcommand{\autor}{Lucas Dias Noronha}
\newcommand{\data}{\today}
\newcommand{\versao}{2.0}

% ---------------------------
% Início do documento
% ---------------------------
\begin{document}

% ==========================================
% PÁGINA 1: CAPA
% ==========================================
\thispagestyle{empty}

\vspace*{\fill}

\begin{center}
    {\Huge \textcolor{AzulCorporativo}{\textbf{Análise de Requisitos}}}\\[0.5cm]
    {\Large \textcolor{AzulPetroleo}{\projeto}}
\end{center}

\vspace*{\fill}

\noindent
\textbf{\textcolor{AzulCorporativo}{Versão:}} \versao
\hfill
\textbf{\textcolor{AzulCorporativo}{Data:}} \data

\newpage

% ==========================================
% PÁGINA 2: HISTÓRICO DE REVISÕES
% ==========================================
\thispagestyle{empty}

\section*{Histórico de Revisões}
\addcontentsline{toc}{section}{Histórico de Revisões}

\vspace{0.5cm}

\begin{center}
\renewcommand{\arraystretch}{1.5}

\begin{tabularx}{\textwidth}{
|>{\centering\arraybackslash}p{1.6cm}
|>{\centering\arraybackslash}p{3.5cm}
|>{\RaggedRight\arraybackslash}X
|>{\RaggedRight\arraybackslash}p{3.5cm}|}
\hline
\rowcolor{AzulCorporativo}
\textcolor{white}{\textbf{Versão}} &
\textcolor{white}{\textbf{Data}} &
\textcolor{white}{\textbf{Descrição das Alterações}} &
\textcolor{white}{\textbf{Autor}} \\ \hline

1.0 & 23 de Março de 2026 &
Criação do documento e definição do problema de negócio e requisitos iniciais. &
\autor \\ \hline

1.1 & 27 de Abril de 2026 &
Definição do escopo e processos contemplados pelo sistema. &
\autor \\ \hline

1.2 & 28 de Abril de 2026 &
Identificação dos stakeholders e definição dos perfis de usuários e interessados no sistema. &
\autor \\ \hline

1.3 & 1 de Maio de 2026 &
Conclusão do documento de requisitos, após revisão de coerência, consolidação das informações e padronização do conteúdo. &
\autor \\ \hline

2.0 & 25 de Agosto de 2026 &
Revisão do escopo do projeto, anteriormente orientado à análise de desempenho logístico, para a construção de um modelo de dados relacional destinado a análises multidimensionais de e-commerce. Revisão do problema de negócio, escopo, requisitos, operações sobre os dados, stakeholders e critérios de aceitação. &
\autor \\ \hline

\end{tabularx}
\end{center}

\newpage

% ==========================================
% PÁGINA 3: SUMÁRIO
% ==========================================
\tableofcontents

\newpage

% ==========================================
% CONTEÚDO TÉCNICO
% ==========================================

\section{Introdução}

Este documento apresenta a análise de requisitos para a construção de um modelo de dados relacional destinado à exploração e análise de dados de comércio eletrônico. O projeto utiliza como fonte o conjunto de dados público \textit{Brazilian E-Commerce Public Dataset by Olist}, composto por informações relacionadas a pedidos, clientes, produtos, vendedores, pagamentos, avaliações e localização geográfica.

O objetivo consiste em estruturar e integrar essas diferentes entidades em um banco de dados relacional consistente, preservando suas relações de negócio e proporcionando uma base adequada para consultas e análises multidimensionais.

A solução proposta não se restringe a uma única perspectiva do comércio eletrônico. A estrutura deverá permitir análises comerciais, comportamentais, geográficas, financeiras e relacionadas ao processo de entrega, conforme as informações disponíveis no conjunto de dados.

A análise de requisitos contempla a definição do problema, a delimitação do domínio, a identificação das entidades e relacionamentos, os requisitos necessários à construção da base analítica e os critérios que orientarão as etapas posteriores de modelagem lógica e física.

% ---------------------------
\section{Problema de Negócio}

\subsection{Definição do Problema}

Os dados de uma operação de comércio eletrônico representam diferentes dimensões do negócio, incluindo clientes, pedidos, produtos, vendedores, pagamentos, avaliações e entregas. Quando essas informações se encontram distribuídas entre diferentes conjuntos de dados, sua utilização para análises integradas exige a compreensão e estruturação adequada das relações existentes entre as entidades.

O problema abordado pelo projeto consiste, portanto, na necessidade de transformar dados relacionais provenientes de diferentes arquivos em uma estrutura de banco de dados consistente e preparada para exploração analítica.

\subsection{Descrição do Problema}

O dataset utilizado no projeto é disponibilizado por meio de múltiplos arquivos, cada um representando diferentes entidades e aspectos do ecossistema de e-commerce. Embora existam identificadores que possibilitem o relacionamento entre esses dados, sua utilização analítica depende da correta definição de entidades, cardinalidades, chaves, restrições e regras de integridade.

A ausência de uma estrutura integrada dificulta consultas que dependem simultaneamente de diferentes dimensões do negócio, como a relação entre clientes e pedidos, produtos e vendedores, pagamentos e vendas, localização geográfica e entregas, ou avaliações e experiência de compra.

O projeto busca solucionar essa necessidade por meio da construção de um modelo relacional capaz de consolidar essas informações e fornecer uma base estruturada para análises posteriores.

\subsection{Impactos}

A ausência de uma estrutura integrada pode resultar em:

\begin{itemize}
    \item dificuldade na realização de análises que envolvam múltiplas entidades;
    \item maior complexidade na preparação dos dados para consultas analíticas;
    \item risco de inconsistências decorrentes de relacionamentos inadequados;
    \item dificuldade na reprodução e validação das análises;
    \item baixa reutilização dos dados entre diferentes perspectivas analíticas;
    \item aumento do esforço necessário para transformar os dados brutos em informações úteis para tomada de decisão.
\end{itemize}

% ---------------------------
\section{Escopo do Sistema}

O projeto contempla a construção de um modelo de dados relacional destinado à integração, organização e exploração analítica dos dados públicos de comércio eletrônico disponibilizados pela Olist.

O modelo deverá representar as principais entidades do domínio e preservar os relacionamentos necessários para permitir análises integradas sobre diferentes dimensões do e-commerce.

Estão incluídos no escopo:

\begin{itemize}
    \item estruturação das entidades e relacionamentos presentes no dataset;
    \item definição de chaves primárias e estrangeiras;
    \item aplicação de regras de integridade referencial;
    \item tratamento e padronização dos dados necessários à carga;
    \item integração dos diferentes conjuntos de dados em banco relacional;
    \item armazenamento de informações relacionadas a clientes, pedidos, itens, produtos, vendedores, pagamentos, avaliações e geolocalização;
    \item suporte a consultas analíticas envolvendo múltiplas entidades;
    \item suporte a análises comerciais, de clientes, produtos, vendedores, pagamentos, avaliações, localização geográfica e entregas;
    \item disponibilização da base para utilização por ferramentas de análise de dados e Business Intelligence.
\end{itemize}

Estão fora do escopo:

\begin{itemize}
    \item desenvolvimento de uma plataforma transacional de e-commerce;
    \item desenvolvimento de funcionalidades de compra ou venda de produtos;
    \item gerenciamento operacional de estoque;
    \item gestão de transportadoras, veículos, motoristas ou rotas;
    \item rastreamento físico de pedidos em tempo real;
    \item processamento de transações financeiras;
    \item atualização dos dados em tempo real;
    \item utilização de informações não disponibilizadas pelo dataset;
    \item utilização de dados pessoais sensíveis ou informações que permitam identificação direta dos consumidores.
\end{itemize}

% ---------------------------
\section{Levantamento de Requisitos}

\subsection{Finalidade do Sistema}

A solução tem por finalidade estruturar, integrar e disponibilizar os dados do ecossistema de e-commerce representado pelo dataset da Olist em um banco de dados relacional preparado para exploração analítica.

A estrutura deverá possibilitar a realização de consultas que relacionem diferentes dimensões do negócio, permitindo investigar características de vendas, clientes, produtos, vendedores, pagamentos, avaliações, distribuição geográfica e processo de entrega.

O banco de dados constitui, portanto, a camada estruturada sobre a qual serão desenvolvidas análises exploratórias, diagnósticas, relatórios, indicadores e demais produtos analíticos.

% ---------------------------
\subsection{Delimitação do Domínio}

A solução abrangerá as entidades Cliente, Pedido, Item do Pedido, Produto, Vendedor, Pagamento, Avaliação e Geolocalização, conforme disponibilizadas no dataset utilizado.

Serão considerados atributos necessários à representação das relações comerciais e analíticas existentes entre essas entidades, incluindo identificadores, informações temporais, valores financeiros, características dos produtos, avaliações e informações geográficas agregadas.

A geolocalização será utilizada para complementar as informações territoriais associadas aos prefixos de CEP de clientes e vendedores, possibilitando análises espaciais e estimativas de distância geográfica.

Ficam explicitamente excluídos do escopo:

\begin{itemize}
    \item dados pessoais diretamente identificáveis;
    \item dados sensíveis;
    \item endereços completos;
    \item informações operacionais de transportadoras;
    \item dados de veículos, motoristas ou rotas efetivamente percorridas;
    \item informações de estoque não disponibilizadas;
    \item quaisquer informações externas não incorporadas formalmente ao projeto.
\end{itemize}

% ---------------------------
\subsection{Entidades Principais}

\begin{itemize}
    \item Cliente
    \item Pedido
    \item Item do Pedido
    \item Produto
    \item Vendedor
    \item Pagamento
    \item Avaliação
    \item Geolocalização
\end{itemize}

% ---------------------------
\clearpage
\mbox{}\vspace{-\baselineskip}
\subsection{Atributos}

\begin{itemize}
    \item \textbf{Cliente:} id\_cliente, id\_cliente\_unico, prefixo\_cep, cidade, estado

    \item \textbf{Pedido:} \nolinkurl{id_pedido}, \nolinkurl{id_cliente},
    \nolinkurl{status}, \nolinkurl{data_compra}, \nolinkurl{data_aprovacao},
    \nolinkurl{data_envio_transportador}, \nolinkurl{data_entrega},
    \nolinkurl{data_estimada}

    \item \textbf{Item do Pedido:} id\_pedido, id\_item, id\_produto, id\_vendedor, data\_limite\_envio, preco, valor\_frete

    \item \textbf{Produto:} id\_produto, categoria, comprimento\_nome, comprimento\_descricao, quantidade\_fotos, peso, comprimento, altura, largura

    \item \textbf{Vendedor:} id\_vendedor, prefixo\_cep, cidade, estado

    \item \textbf{Pagamento:} id\_pedido, sequencial\_pagamento, tipo\_pagamento, parcelas, valor

    \item \textbf{Avaliação:} id\_avaliacao, id\_pedido, nota, data\_criacao, data\_resposta

    \item \textbf{Geolocalização:} prefixo\_cep, latitude, longitude, cidade, estado
\end{itemize}

% ---------------------------
\subsection{Relacionamentos}

\textbf{Entidades}: Cliente, Pedido, Item do Pedido, Produto, Vendedor, Pagamento, Avaliação, Geolocalização

\textbf{Relacionamentos:}

\begin{itemize}
    \item Cliente realiza Pedido (1:N)
    \item Pedido possui Item do Pedido (1:N)
    \item Produto compõe Item do Pedido (1:N)
    \item Vendedor fornece Item do Pedido (1:N)
    \item Pedido gera Pagamento (1:N)
    \item Pedido pode receber Avaliação (1:1)
    \item Cliente pode ser associado a registros de Geolocalização por meio do prefixo de CEP
    \item Vendedor pode ser associado a registros de Geolocalização por meio do prefixo de CEP
\end{itemize}

% ---------------------------
\clearpage
\subsection{Requisitos Funcionais}

Os requisitos funcionais descrevem as capacidades que a estrutura de dados deverá oferecer para possibilitar a integração e exploração analítica das informações.

\renewcommand{\arraystretch}{1.15}
\begin{longtable}{|>{\RaggedRight\arraybackslash}p{2cm}|>{\RaggedRight\arraybackslash}p{12cm}|}
    \hline
    \rowcolor{AzulCorporativo}
    \textcolor{white}{\textbf{ID}} &
    \textcolor{white}{\textbf{Descrição}} \\ \hline
    \endfirsthead

    \hline
    \rowcolor{AzulCorporativo}
    \textcolor{white}{\textbf{ID}} &
    \textcolor{white}{\textbf{Descrição}} \\ \hline
    \endhead

    RF01 & O modelo deve permitir a integração dos dados de clientes, pedidos, itens, produtos, vendedores, pagamentos, avaliações e geolocalização. \\ \hline

    RF02 & O modelo deve permitir a recuperação do histórico de pedidos e suas respectivas informações relacionadas. \\ \hline

    RF03 & O modelo deve permitir a análise de vendas segundo diferentes dimensões, incluindo período, produto, categoria, cliente, vendedor e localização geográfica. \\ \hline

    RF04 & O modelo deve permitir a análise do comportamento de compra dos clientes a partir do histórico disponível. \\ \hline

    RF05 & O modelo deve permitir a análise do desempenho comercial de produtos, categorias e vendedores. \\ \hline

    RF06 & O modelo deve permitir a análise das formas, valores e condições de pagamento associadas aos pedidos. \\ \hline

    RF07 & O modelo deve permitir relacionar avaliações dos clientes às demais informações dos respectivos pedidos. \\ \hline

    RF08 & O modelo deve permitir analisar datas previstas e efetivas associadas ao processamento e à entrega dos pedidos. \\ \hline

    RF09 & O modelo deve permitir análises geográficas utilizando informações de clientes e vendedores e, quando aplicável, coordenadas associadas aos prefixos de CEP. \\ \hline

    RF10 & O modelo deve permitir consultas que combinem múltiplas entidades sem comprometer a integridade dos relacionamentos. \\ \hline

    RF11 & O modelo deve permitir agregações e filtros por dimensões temporais, comerciais, geográficas e relacionadas aos pedidos. \\ \hline

    RF12 & O modelo deve fornecer uma base estruturada para consumo por ferramentas de análise de dados e Business Intelligence. \\ \hline
\end{longtable}

% ---------------------------
\subsection{Requisitos Não Funcionais}

Os requisitos não funcionais descrevem restrições e qualidades que o sistema deve atender.

\begin{longtable}{|>{\RaggedRight\arraybackslash}p{2cm}|>{\RaggedRight\arraybackslash}p{12cm}|}
    \hline
    \rowcolor{AzulCorporativo}
    \textcolor{white}{\textbf{ID}} &
    \textcolor{white}{\textbf{Descrição}} \\ \hline
    \endfirsthead

    \hline
    \rowcolor{AzulCorporativo}
    \textcolor{white}{\textbf{ID}} &
    \textcolor{white}{\textbf{Descrição}} \\ \hline
    \endhead

    RNF01 & O sistema deve garantir a integridade e consistência dos dados armazenados. \\ \hline

    RNF02 & O sistema deve assegurar a correta integração entre as diferentes entidades do modelo de dados. \\ \hline

    RNF03 & O sistema deve permitir consultas com desempenho adequado para análise de dados. \\ \hline

    RNF04 & O sistema deve ser estruturado de forma a facilitar escalabilidade, manutenção e extensão. \\ \hline

    RNF05 & O sistema deve garantir padronização nos formatos de dados, especialmente datas, identificadores e atributos geográficos. \\ \hline

    RNF06 & O sistema deve assegurar a rastreabilidade das informações ao longo do fluxo de ingestão e transformação dos dados. \\ \hline

    RNF07 & O sistema deve ser projetado de forma a permitir fácil extensão para novas análises. \\ \hline

    RNF08 & O sistema deve incorporar verificações de qualidade dos dados, incluindo duplicidades, valores ausentes e inconsistências de relacionamento. \\ \hline
\end{longtable}

% ---------------------------
\section{Regras de Negócio}

\begin{itemize}
    \item Um cliente pode realizar múltiplos pedidos.
    \item Cada pedido deve estar associado a um único cliente.
    \item Um pedido deve possuir pelo menos um item quando houver registro correspondente na tabela de itens.
    \item Cada item de pedido deve estar associado a um único produto.
    \item Cada item de pedido deve estar associado a um único vendedor.
    \item Um pedido pode possuir um ou mais registros de pagamento.
    \item Um pedido pode possuir uma avaliação associada.
    \item Um pedido não pode existir sem um cliente associado.
    \item Um mesmo pedido pode conter itens fornecidos por vendedores distintos.
    \item Informações geográficas de clientes e vendedores devem ser associadas por meio do prefixo de CEP, observadas as particularidades do dataset de geolocalização.
    \item As análises devem respeitar a granularidade de cada entidade para evitar duplicações decorrentes de junções entre tabelas com cardinalidades diferentes.
\end{itemize}

% ---------------------------
\section{Operações sobre os Dados}

Considerando a natureza analítica do projeto, as operações sobre os dados não correspondem ao fluxo CRUD de uma aplicação transacional. O banco será construído a partir de dados históricos provenientes do dataset utilizado e terá como finalidade principal a integração, persistência e consulta das informações.

\subsection{Extração}

Os dados serão obtidos a partir dos arquivos que compõem o dataset original, preservando-se uma camada de dados brutos para rastreabilidade.

\subsection{Transformação}

Os dados poderão passar por operações de:

\begin{itemize}
    \item padronização de tipos;
    \item tratamento de valores ausentes;
    \item validação de identificadores;
    \item adequação de datas;
    \item verificação de duplicidades;
    \item validação de relacionamentos;
    \item preparação dos dados para carga no banco relacional.
\end{itemize}

\subsection{Carga}

Os registros tratados serão carregados nas respectivas tabelas do banco de dados, respeitando dependências, chaves e restrições de integridade referencial.

\subsection{Consulta}

A principal forma de utilização da base será por meio de consultas analíticas, incluindo:

\begin{itemize}
    \item consultas envolvendo múltiplas entidades;
    \item agregações temporais;
    \item análises de vendas;
    \item análises de clientes;
    \item análises de produtos e categorias;
    \item análises de vendedores;
    \item análises de pagamentos;
    \item análises de avaliações;
    \item análises geográficas;
    \item análises relacionadas ao processamento e à entrega dos pedidos.
\end{itemize}

% ---------------------------
\section{Modelo de Dados (visão preliminar)}

\subsection{Tipo de Banco de Dados}

Considerando a natureza estruturada dos dados e a necessidade de manter relações consistentes entre entidades, o sistema será baseado em um modelo de dados relacional.

A implementação física deverá utilizar um Sistema Gerenciador de Banco de Dados relacional compatível com os requisitos do projeto, sendo o PostgreSQL a tecnologia adotada para esta implementação.

\subsection{Decisões Iniciais de Modelagem}

\begin{itemize}
    \item \textbf{Normalização:} o modelo será estruturado de forma normalizada, buscando reduzir redundâncias e preservar a consistência entre as entidades.

    \item \textbf{Índices:} serão criados índices em atributos frequentemente utilizados em operações de junção e filtragem, como identificadores primários e estrangeiros.

    \item \textbf{Integridade Referencial:} serão definidas chaves primárias e estrangeiras para assegurar a consistência entre as entidades relacionadas.

    \item \textbf{Granularidade:} as tabelas deverão preservar a granularidade original das entidades, evitando consolidações que comprometam análises futuras.

    \item \textbf{Escalabilidade:} a estrutura deverá permitir futuras extensões e inclusão de novas camadas analíticas sem necessidade de reconstrução integral do modelo.
\end{itemize}

% ---------------------------
\section{Interfaces e Integrações}

A estrutura deverá possibilitar o consumo dos dados por ferramentas utilizadas nas etapas de exploração, análise e visualização.

\begin{itemize}
    \item \textbf{SQL:} acesso ao banco relacional para realização de consultas, agregações e análises envolvendo as entidades modeladas.

    \item \textbf{Python:} integração com bibliotecas de análise de dados para exploração, transformação complementar, validação e desenvolvimento de análises.

    \item \textbf{Ferramentas de Business Intelligence:} possibilidade de conexão com ferramentas de visualização para construção de dashboards e relatórios analíticos.

    \item \textbf{Arquivos estruturados:} utilização dos arquivos originais como fonte de dados e possibilidade de geração de conjuntos derivados para etapas específicas da análise.
\end{itemize}

% ---------------------------
\section{Restrições e Premissas}

\subsection{Restrições}

\begin{itemize}
    \item O sistema será limitado aos dados disponibilizados pelo dataset utilizado, salvo incorporação formal de novas fontes em versão futura.

    \item Não serão considerados dados pessoais sensíveis ou diretamente identificáveis.

    \item O sistema não contemplará, nesta versão, dados não estruturados extensos que não façam parte do escopo definido.

    \item A modelagem será baseada em dados históricos, não contemplando atualização em tempo real.

    \item Não serão realizadas análises operacionais de transportadoras, veículos, motoristas ou rotas, pois tais informações não estão presentes no dataset.

    \item A geolocalização representa posições associadas a prefixos de CEP e não trajetos efetivamente percorridos pelos pedidos.
\end{itemize}

\subsection{Premissas}

\begin{itemize}
    \item Assume-se que o dataset disponibilizado pela Olist constitui a fonte primária de dados do projeto.

    \item A qualidade, completude e consistência dos dados não serão presumidas, devendo ser verificadas durante as etapas de exploração, tratamento e carga.

    \item Os identificadores disponibilizados serão avaliados quanto à unicidade e integridade antes de sua utilização como chaves no modelo.

    \item O sistema será utilizado predominantemente para fins analíticos e de apoio à tomada de decisão.

    \item O volume de dados é compatível com processamento e armazenamento em um Sistema Gerenciador de Banco de Dados relacional convencional.

    \item As análises estarão limitadas às informações efetivamente representadas no dataset e às variáveis que possam ser derivadas delas de maneira metodologicamente justificável.
\end{itemize}

% ---------------------------
\section{Critérios de Aceitação}

O modelo será considerado adequado quando atender aos seguintes critérios:

\begin{itemize}
    \item As entidades definidas no escopo devem estar implementadas e relacionadas conforme o modelo lógico aprovado.

    \item As chaves primárias e estrangeiras devem preservar a integridade referencial entre as entidades.

    \item O processo de carga deve ser capaz de incorporar os dados necessários do dataset ao banco relacional.

    \item Os registros carregados devem passar por verificações de qualidade, consistência, tipos e relacionamentos.

    \item O banco deve permitir consultas envolvendo múltiplas entidades sem gerar relacionamentos incorretos ou duplicações decorrentes de junções inadequadas.

    \item Deve ser possível realizar análises relacionadas a clientes, pedidos, vendas, produtos, vendedores, pagamentos, avaliações, localização geográfica e entregas.

    \item O modelo deve ser capaz de processar o volume completo de dados definido para o projeto.

    \item As principais consultas analíticas devem apresentar desempenho adequado ao volume de dados utilizado.

    \item A estrutura deve possuir documentação suficiente para permitir a compreensão das entidades, atributos, relacionamentos e regras de negócio.
\end{itemize}

% ---------------------------
\section{Dicionário de Dados}

Abaixo está o detalhamento dos principais campos considerados para a modelagem do banco de dados, com base no dataset utilizado.

\vspace{0.5cm}

{\footnotesize
\begin{longtable}{
    |>{\raggedright\arraybackslash}p{5.5cm}
    |>{\raggedright\arraybackslash}p{2cm}
    |>{\raggedright\arraybackslash}p{6.8cm}|
}
\hline
\rowcolor{AzulCorporativo}
\textcolor{white}{\textbf{Campo}} &
\textcolor{white}{\textbf{Tipo}} &
\textcolor{white}{\textbf{Descrição}} \\ \hline
\endfirsthead

\hline
\rowcolor{AzulCorporativo}
\textcolor{white}{\textbf{Campo}} &
\textcolor{white}{\textbf{Tipo}} &
\textcolor{white}{\textbf{Descrição}} \\ \hline
\endhead

customer\_id & VARCHAR & Identificador do cliente associado ao pedido \\
customer\_unique\_id & VARCHAR & Identificador persistente do cliente no dataset \\
customer\_zip\_code\_prefix & INT & Prefixo do CEP do cliente \\
customer\_city & VARCHAR & Cidade do cliente \\
customer\_state & VARCHAR & Estado do cliente \\

order\_id & VARCHAR & Identificador único do pedido \\
order\_status & VARCHAR & Status do pedido \\
order\_purchase\_timestamp & TIMESTAMP & Data e hora da compra \\
order\_approved\_at & TIMESTAMP & Data e hora de aprovação do pedido \\
order\_delivered\_carrier\_date & TIMESTAMP & Data de entrega do pedido ao transportador \\
order\_delivered\_customer\_date & TIMESTAMP & Data de entrega do pedido ao cliente \\
order\_estimated\_delivery\_date & TIMESTAMP & Data estimada de entrega \\

order\_item\_id & INT & Número sequencial do item dentro do pedido \\
shipping\_limit\_date & TIMESTAMP & Data limite de envio do item pelo vendedor \\
price & NUMERIC & Preço do item \\
freight\_value & NUMERIC & Valor de frete associado ao item \\

product\_id & VARCHAR & Identificador único do produto \\
product\_category\_name & VARCHAR & Categoria do produto \\
product\_name\_length & INT & Comprimento do nome do produto \\
product\_description\_length & INT & Comprimento da descrição do produto \\
product\_photos\_qty & INT & Quantidade de fotos associadas ao produto \\
product\_weight\_g & INT & Peso do produto em gramas \\
product\_length\_cm & INT & Comprimento do produto em centímetros \\
product\_height\_cm & INT & Altura do produto em centímetros \\
product\_width\_cm & INT & Largura do produto em centímetros \\

seller\_id & VARCHAR & Identificador do vendedor \\
seller\_zip\_code\_prefix & INT & Prefixo do CEP do vendedor \\
seller\_city & VARCHAR & Cidade do vendedor \\
seller\_state & VARCHAR & Estado do vendedor \\

payment\_sequential & INT & Sequência do pagamento no pedido \\
payment\_type & VARCHAR & Tipo de pagamento \\
payment\_installments & INT & Número de parcelas \\
payment\_value & NUMERIC & Valor do pagamento \\

review\_id & VARCHAR & Identificador da avaliação \\
review\_score & INT & Nota da avaliação do pedido \\
review\_creation\_date & TIMESTAMP & Data de criação da avaliação \\
review\_answer\_timestamp & TIMESTAMP & Data e hora da resposta à avaliação \\

geolocation\_zip\_code\_prefix & INT & Prefixo de CEP utilizado na base de geolocalização \\
geolocation\_lat & NUMERIC & Latitude associada ao prefixo de CEP \\
geolocation\_lng & NUMERIC & Longitude associada ao prefixo de CEP \\
geolocation\_city & VARCHAR & Cidade associada ao registro de geolocalização \\
geolocation\_state & VARCHAR & Estado associado ao registro de geolocalização \\ \hline

\end{longtable}
}

% ---------------------------
\section{Identificação de Stakeholders}

Esta seção identifica os principais perfis interessados na estruturação e utilização das informações disponibilizadas pelo modelo analítico.

\vspace{0.5cm}

\renewcommand{\arraystretch}{1.15}
\begin{longtable}{
|>{\RaggedRight\arraybackslash}p{3.7cm}
|>{\RaggedRight\arraybackslash}p{3.7cm}
|>{\RaggedRight\arraybackslash}p{6.9cm}|}
    \hline
    \rowcolor{AzulCorporativo}
    \textcolor{white}{\textbf{Stakeholder}} &
    \textcolor{white}{\textbf{Papel}} &
    \textcolor{white}{\textbf{Interesse no Sistema}} \\ \hline
    \endfirsthead

    \hline
    \rowcolor{AzulCorporativo}
    \textcolor{white}{\textbf{Stakeholder}} &
    \textcolor{white}{\textbf{Papel}} &
    \textcolor{white}{\textbf{Interesse no Sistema}} \\ \hline
    \endhead

    Gestores de Negócio &
    Tomadores de decisão &
    Utilizar informações consolidadas para compreender desempenho comercial, comportamento dos clientes e demais dimensões relevantes do negócio. \\ \hline

    Analistas de Dados &
    Usuários analíticos &
    Realizar consultas, análises exploratórias e diagnósticas e produzir indicadores e informações para apoio à decisão. \\ \hline

    Profissionais de BI &
    Consumidores da camada de dados &
    Utilizar dados estruturados para construção de relatórios, indicadores e dashboards. \\ \hline

    Engenharia de Dados / Administração de Banco de Dados &
    Responsáveis pela estrutura de dados &
    Assegurar qualidade, integridade, organização, desempenho e manutenção da estrutura relacional. \\ \hline

    Áreas de Negócio &
    Consumidores das análises &
    Utilizar os resultados das análises para compreender clientes, vendas, produtos, vendedores, pagamentos, experiência de compra e entregas. \\ \hline
\end{longtable}

% ---------------------------
\section{Considerações Finais}

A análise de requisitos apresentada neste documento estabelece os fundamentos para a construção de um modelo de dados relacional destinado à integração e exploração analítica de informações de comércio eletrônico.

A estrutura proposta busca representar de maneira consistente as diferentes entidades presentes no dataset da Olist, preservando seus relacionamentos e permitindo que informações de clientes, pedidos, produtos, vendedores, pagamentos, avaliações, geolocalização e entregas sejam analisadas de maneira integrada.

Diferentemente de uma solução orientada exclusivamente ao desempenho logístico, o modelo possui caráter multidimensional. As informações relacionadas às entregas constituem uma das perspectivas possíveis de análise, juntamente com dimensões comerciais, comportamentais, financeiras, geográficas e relacionadas à experiência do cliente.

A definição das regras de negócio, requisitos, restrições e critérios de aceitação contribui para assegurar coerência entre a estrutura relacional e sua finalidade analítica, além de estabelecer limites claros quanto às conclusões que podem ser obtidas a partir dos dados disponíveis.

Dessa forma, o documento consolida as diretrizes que orientarão as etapas subsequentes de modelagem lógica, implementação física, carga dos dados e desenvolvimento das análises.

% ---------------------------
\clearpage
\section{Referências}

\begin{itemize}

\item \textbf{Dataset Olist E-commerce} \\
Disponível em: \url{https://www.kaggle.com/datasets/olistbr/brazilian-ecommerce} \\
Base pública contendo informações de pedidos, clientes, vendedores, pagamentos, avaliações, produtos e geolocalização no contexto do comércio eletrônico brasileiro.

\item \textbf{Documentação do PostgreSQL} \\
Disponível em: \url{https://www.postgresql.org/docs/} \\
Referência oficial para modelagem relacional, tipos de dados, integridade referencial e implementação em banco de dados.

\item \textbf{Documentação do DBeaver} \\
Disponível em: \url{https://dbeaver.io/documentation/} \\
Ferramenta utilizada para gerenciamento, consulta e visualização do banco de dados.

\item \textbf{Elmasri, R.; Navathe, S.} \\
\textit{Fundamentals of Database Systems}. \\
Livro de referência para conceitos de modelagem de dados, normalização e projeto de bancos relacionais.

\item \textbf{Silberschatz, A.; Korth, H.; Sudarshan, S.} \\
\textit{Database System Concepts}. \\
Obra clássica que fundamenta princípios teóricos de sistemas de banco de dados.

\item \textbf{Date, C. J.} \\
\textit{An Introduction to Database Systems}. \\
Referência consolidada sobre teoria relacional e boas práticas em modelagem de dados.

\end{itemize}

\end{document}
